\ChapterImageStar[cap:resultados]{Resultados}{./images/fondo.png}\label{cap:resultados}
\mbox{}\\

\section{Caracterización del Grupo~\GRID~e identificación de necesidades}
\noindent

La caracterización del Grupo de Investigación en Redes, Información y Distribución (\GRID) se desarrolló con el objetivo de contextualizar la propuesta de la notación NDTI dentro del entorno académico en el cual será aplicada, así como identificar las necesidades específicas relacionadas con el modelado de arquitecturas de TI en el programa de Ingeniería de Sistemas y Computación.

El Grupo~\GRID~orienta su actividad investigativa hacia áreas como redes de computadores, sistemas operativos, sistemas distribuidos y seguridad informática. Estas áreas implican el análisis y diseño de sistemas complejos que abarcan múltiples niveles de abstracción, desde la infraestructura física hasta la lógica de aplicaciones y los procesos organizacionales. En consecuencia, el modelado de estos sistemas constituye una actividad fundamental tanto en el ámbito investigativo como en el formativo.

A pesar de la relevancia del modelado dentro de estas áreas, no existe actualmente una notación unificada que permita representar de manera coherente los distintos componentes de la infraestructura de TI dentro del contexto académico del programa. Esta situación genera una fragmentación en la forma en que se conceptualizan y comunican las arquitecturas, dificultando la comprensión integral de los sistemas por parte de los estudiantes.

Con el fin de identificar de manera precisa estas problemáticas, se realizó un proceso de recolección de información con los integrantes del Grupo \GRID, enfocado en analizar su relación con el currículo, las herramientas de modelado utilizadas, sus preferencias, y las dificultades asociadas a la representación de infraestructuras tecnológicas. Este análisis permitió evidenciar no solo la diversidad de enfoques existentes, sino también la necesidad de una herramienta que articule de forma clara los distintos niveles de la arquitectura.

A partir de esta caracterización, se establecen los requerimientos funcionales y conceptuales que fundamentan el diseño de la notación NDTI, orientada a ofrecer una solución que combine claridad semántica, capacidad de integración multicapa y aplicabilidad en contextos educativos.

En las siguientes subsecciones se presentan los resultados detallados de este proceso de caracterización, organizados en torno a diferentes dimensiones de análisis.

\section{Revisión sistemática de la literatura}
\noindent

Con el propósito de fundamentar el desarrollo de la notación NDTI y garantizar decisiones de diseño respaldadas en evidencia, se llevó a cabo un \textbf{estudio de mapeo sistemático (SMS)} orientado a identificar, clasificar y analizar trabajos relacionados con notaciones, metodologías y herramientas de modelado en el contexto de la infraestructura de TI.

Este estudio se desarrolló siguiendo un enfoque metodológico estructurado, basado en el modelo Goal-Question-Metric (GQM), el cual permitió definir de manera explícita las metas del análisis, las preguntas de investigación y las métricas utilizadas para evaluar los resultados. A partir de esta base, se diseñó una estrategia de búsqueda combinada que integró consultas en bases de datos académicas reconocidas (ACM, IEEE Xplore, Springer, ScienceDirect y Taylor \& Francis), junto con la técnica de búsqueda por bola de nieve, con el fin de ampliar la cobertura y garantizar la inclusión de estudios relevantes.

Como resultado de la aplicación de las cadenas de búsqueda, se identificó un conjunto inicial de 13\,376 registros. Tras la aplicación de criterios de inclusión y exclusión, este conjunto se redujo a 559 estudios potencialmente relevantes. Posteriormente, la eliminación de duplicados y el proceso de priorización basado en métricas de calidad permitieron identificar 68 estudios directamente alineados con los objetivos de la investigación. Finalmente, mediante la aplicación de la estrategia de bola de nieve, se incorporaron 63 estudios adicionales, consolidando un corpus final de 131 artículos para su análisis.

El análisis de los estudios seleccionados permitió identificar tendencias relevantes en el uso de metodologías, herramientas y enfoques de modelado en el ámbito de la infraestructura de TI. Entre los principales hallazgos se destaca la predominancia de enfoques centrados en modelos, metodologías y notaciones, así como la presencia de áreas relacionadas como HPC y HTC, las cuales evidencian la diversidad de contextos en los que se aplican soluciones de modelado.

Asimismo, se observó una alta cobertura de las preguntas de investigación planteadas, donde el 90.07\% de los estudios abordaron aspectos relacionados con metodologías de modelado (RQ1), mientras que el 99.23\% evidenciaron el uso o análisis de herramientas tecnológicas asociadas (RQ2). Estos resultados reflejan un interés significativo en la literatura por el desarrollo de soluciones de modelado, aunque también ponen en evidencia la falta de estandarización y la fragmentación de enfoques en el diseño de infraestructura de TI.

En conjunto, los resultados del SMS evidencian la necesidad de contar con una notación que integre múltiples dimensiones del modelado de infraestructura, proporcionando coherencia semántica, capacidad de integración entre capas y facilidad de uso en contextos académicos y prácticos. Esta necesidad constituye el principal fundamento para el desarrollo de la notación NDTI propuesta en este trabajo.

\section{Análisis de Decisiones y Resolución (DAR)}
\noindent

La aplicación de la metodología DAR del modelo CMMI permitió evaluar de manera objetiva distintas notaciones de modelado arquitectónico mediante un conjunto de criterios técnicos, metodológicos y pedagógicos previamente definidos. Entre los criterios considerados se incluyeron la popularidad de la notación, la disponibilidad de documentación, el nivel de formalidad y respaldo institucional, la vigencia y el grado de minimalismo en su representación.

El proceso de evaluación sistemática evidenció que las notaciones \textbf{UML}, \textbf{BPMN} y \textbf{ArchiMate} presentan un desempeño similar en términos globales, obteniendo puntuaciones cercanas sin que exista una diferencia suficientemente significativa que permita seleccionar una única alternativa como dominante.

Ante esta situación, se adoptó una decisión de diseño basada en la combinación de enfoques, seleccionando a \textit{ArchiMate} como eje estructural de la propuesta debido a su madurez como estándar, su capacidad de modelado multicapa y su cobertura de diferentes niveles de abstracción. Esta base se complementa con principios y elementos provenientes de \textit{UML} y \textit{BPMN}, con el fin de fortalecer aspectos específicos relacionados con la claridad estructural y la representación de procesos.

Los criterios que favorecieron esta decisión incluyen la amplia disponibilidad de documentación, el respaldo institucional de las notaciones evaluadas, su adopción en contextos académicos y profesionales, así como la complementariedad entre sus capacidades de modelado. De esta manera, la propuesta \NDITI~se fundamenta en una integración controlada de enfoques, orientada a maximizar su aplicabilidad en el contexto académico y técnico del proyecto.

\section{Resultados de la evaluación}
\noindent
En esta sección se presentan los resultados obtenidos a partir de la aplicación del ejercicio de diagramación y la encuesta de percepción, organizados en dos dimensiones: (i) desempeño en la actividad de modelado y (ii) percepción de los participantes sobre la notación NDTI.

\subsection{Resultados de la actividad de diagramación}
\noindent
El análisis de los diagramas se realizó de forma comparativa entre el grupo experimental (uso de la notación NDTI) y el grupo de referencia (diagramación libre), con base en los criterios definidos en la sección metodológica.

\subsubsection{Resumen general de resultados}
\noindent
En la Tabla~\ref{tab:resumen_diagramacion} se presenta el resumen de los resultados obtenidos por ambos grupos en los criterios de evaluación.

\begin{table}[H]
\centering
\caption{Resumen comparativo de resultados en la actividad de diagramación}
\label{tab:resumen_diagramacion}
\begin{tabular}{|l|c|c|}
\hline
\textbf{Criterio} & \textbf{Grupo NDTI} & \textbf{Grupo libre} \\
\hline
Correspondencia con el enunciado & [X.X] & [X.X] \\
Estructuración del modelo & [X.X] & [X.X] \\
Representación de relaciones & [X.X] & [X.X] \\
Completitud del modelo & [X.X] & [X.X] \\
\hline
\end{tabular}
\end{table}

\noindent
Adicionalmente, estos resultados serán representados mediante un \textbf{gráfico de barras comparativo}, con el fin de visualizar las diferencias de desempeño entre ambos grupos por criterio evaluado.

\FloatBarrier

\subsubsection{Análisis cualitativo de los diagramas}
\noindent
A partir de la revisión de los diagramas elaborados por los participantes, se identificaron patrones relevantes en la forma en que ambos grupos abordaron la representación del problema.

\paragraph{Grupo con notación NDTI}
\noindent
[Describir aquí los patrones observados, por ejemplo:
- Uso consistente de capas
- Mejores relaciones
- Uso adecuado de elementos
- Errores recurrentes en conectores específicos]

\paragraph{Grupo sin notación}
\noindent
[Describir:
- Representaciones informales
- Falta de estructura
- Ambigüedad en relaciones
- Inconsistencias en la organización del modelo]

\paragraph{Comparación entre grupos}
\noindent
[Interpretación directa:
- Diferencias clave
- Ventajas de la notación
- Limitaciones observadas]

\subsection{Resultados de la encuesta de percepción}
\noindent
La encuesta fue aplicada al grupo que utilizó la notación NDTI, con el fin de evaluar su experiencia en relación con los ejes definidos.

\subsubsection{Resultados cuantitativos}
\noindent
En la Tabla~\ref{tab:resultados_encuesta} se presentan los promedios obtenidos por cada eje de evaluación.

\begin{table}[H]
\centering
\caption{Resultados de la encuesta de percepción}
\label{tab:resultados_encuesta}
\begin{tabular}{|l|c|}
\hline
\textbf{Eje de evaluación} & \textbf{Promedio (1-5)} \\
\hline
Comprensión de la notación & [X.X] \\
Uso de los elementos & [X.X] \\
Expresividad y utilidad & [X.X] \\
Experiencia general & [X.X] \\
\hline
\end{tabular}
\end{table}

\FloatBarrier

\noindent
Los resultados anteriores serán complementados mediante representaciones gráficas que permitan un análisis más detallado de la distribución de respuestas:

\begin{itemize}
    \item \textbf{Gráficos de barras tipo Likert} para cada afirmación evaluada, mostrando la distribución de respuestas en la escala de 1 a 5.
    \item \textbf{Gráficos de barras agrupadas} por eje de evaluación, permitiendo comparar tendencias entre dimensiones (comprensión, uso, utilidad y experiencia).
    \item \textbf{Diagramas de radar (opcional)} para visualizar de manera integrada el comportamiento de los diferentes ejes de evaluación.
\end{itemize}

\noindent
\textit{Nota: Los valores y gráficos serán completados una vez se procesen las respuestas de la encuesta.}

\subsubsection{Análisis cualitativo de respuestas abiertas}
\noindent
Las respuestas abiertas de la encuesta fueron analizadas mediante categorización temática, agrupando las opiniones de los participantes en categorías representativas.

\begin{table}[H]
\centering
\caption{Síntesis de resultados cualitativos}
\begin{tabular}{|p{4cm}|p{10cm}|}
\hline
\textbf{Categoría} & \textbf{Descripción} \\
\hline
Fortalezas & [Ej: claridad de capas, facilidad de interpretación, organización visual] \\
\hline
Dificultades & [Ej: confusión en algunos elementos, curva de aprendizaje inicial] \\
\hline
Oportunidades de mejora & [Ej: simplificación de símbolos, guías de uso, ejemplos] \\
\hline
\end{tabular}
\end{table}

\noindent
Estas categorías podrán complementarse con un \textbf{gráfico de frecuencias} que indique la recurrencia de cada tipo de comentario.

\subsection{Integración de resultados}
\noindent
La integración de los resultados permitió contrastar el desempeño observado en la actividad de diagramación con la percepción de los participantes sobre la notación NDTI.

\noindent
[Ejemplo de lo que irá aquí cuando tengas datos:]
Se analizará si un mayor nivel de comprensión percibida se traduce en un mejor desempeño en la construcción de diagramas, así como identificar posibles discrepancias entre percepción y aplicación práctica.

\noindent
Asimismo, se buscará identificar relaciones entre:
\begin{itemize}
    \item Nivel de comprensión percibida y calidad del modelo construido.
    \item Facilidad de uso y estructuración del diagrama.
    \item Utilidad percibida y completitud del modelo.
\end{itemize}

\noindent
\textit{Nota: La interpretación profunda de estos hallazgos se desarrollará en la sección de conclusiones.}